% GENERATED FILE. Edit canonical lessons, then run npm run generate:books.
% canonical-source-hash: fnv1a64:d31fde075a494642
% canonical-lessons: MR-C128-cauk, MR-C128-bag, MR-C128-udyan, MR-C128-vimantal, MR-C128-rugnalay

\chapter{Square, Garden, Park, Airport, Hospital}
\label{ch:mr-square-garden-park-airport-hospital}
\hlchaptermodality{128}

\begin{chapteropening}
\textbf{By the end of this chapter:} \emph{I can say a square, a garden, a park, an airport and a hospital.}

Complete the last lesson of chapter 128: I can say a square, a garden, a park, an airport and a hospital.
\end{chapteropening}

\section[\texorpdfstring{\mr{चौक} (cauk)}{चौक (cauk)}]{\mr{चौक} (cauk) — a square, a crossroads}

\label{lesson:MR-C128-cauk}



\begin{quote}
Before the new one: say the Marathi for a turn, then the Marathi for a corner.

Name near and far: \textbf{\mr{हा}} for the one here, \textbf{\mr{तो}} for the one there, and \textbf{\mr{तो}} again for \emph{he}.
\end{quote}

\subsection*{You'll want to know: \mr{चौक}}
\textbf{\mr{चौक}} — \emph{cauk} — \textquotedblleft{}a square, a crossroads\textquotedblright{}.

\begin{cousinweb}
\textbf{\mr{चौक}} (\emph{cauk}) is a tadbhava, a word inherited from Sanskrit and reshaped by sound change. It comes from \emph{catuṣka}, \textquotedblleft{}set of four; quadrangular courtyard\textquotedblright{} (CDIAL 4629). It is built on \emph{catur}, \textquotedblleft{}four\textquotedblright{}.
\end{cousinweb}

\begin{grammarlens}[title={What you've built}]
One more word for this chapter.
\end{grammarlens}

\subsection*{Your turn}
\begin{itemize}
\raggedright
  \item \emph{Say it:} \emph{cauk}
  \item \emph{Say it:} \emph{cauk}, once more
  \item \emph{Say it:} say \emph{cauk}
  \item \emph{Recall:} say the Marathi for a turn, then the Marathi for a corner, then say \emph{cauk} again
\end{itemize}

\begin{tcolorbox}[breakable,colback=teal!4,colframe=teal!35!black,title={Before you move on}]
What does \mr{चौक} mean? (\textquotedblleft{}A square, a crossroads\textquotedblright{}.) Say it once more.
\end{tcolorbox}
\section[\texorpdfstring{\mr{बाग} (bāg)}{बाग (bāg)}]{\mr{बाग} (bāg) — a garden}

\label{lesson:MR-C128-bag}



\begin{quote}
Before the new one: say the Marathi for a corner, then the Marathi for a square.
\end{quote}

\subsection*{You'll want to know: \mr{बाग}}
\textbf{\mr{बाग}} — \emph{bāg} — \textquotedblleft{}a garden\textquotedblright{}.

\begin{cousinweb}
\textbf{\mr{बाग}} (\emph{bāg}) is a loan, usually traced to Persian \emph{bāgh}, \textquotedblleft{}garden\textquotedblright{}.
\end{cousinweb}

\begin{grammarlens}[title={What you've built}]
One more word for this chapter.
\end{grammarlens}

\subsection*{Your turn}
\begin{itemize}
\raggedright
  \item \emph{Say it:} \emph{bāg}
  \item \emph{Say it:} \emph{bāg}, once more
  \item \emph{Say it:} say \emph{bāg}
  \item \emph{Recall:} say the Marathi for a corner, then the Marathi for a square, then say \emph{bāg} again
\end{itemize}

\begin{tcolorbox}[breakable,colback=teal!4,colframe=teal!35!black,title={Before you move on}]
What does \mr{बाग} mean? (\textquotedblleft{}A garden\textquotedblright{}.) Say it once more.
\end{tcolorbox}
\section[\texorpdfstring{\mr{उद्यान} (udyān)}{उद्यान (udyān)}]{\mr{उद्यान} (udyān) — a park}

\label{lesson:MR-C128-udyan}



\begin{quote}
Before the new one: say the Marathi for a square, then the Marathi for a garden.
\end{quote}

\subsection*{You'll want to know: \mr{उद्यान}}
\textbf{\mr{उद्यान}} — \emph{udyān} — \textquotedblleft{}a park\textquotedblright{}.

\begin{cousinweb}
\textbf{\mr{उद्यान}} (\emph{udyān}) is Sanskrit \emph{udyāna}, \textquotedblleft{}going out; park, garden\textquotedblright{}, taken over whole as a learned loan, a tatsama (CDIAL 2052).
\end{cousinweb}

\begin{grammarlens}[title={What you've built}]
One more word for this chapter.
\end{grammarlens}

\subsection*{Your turn}
\begin{itemize}
\raggedright
  \item \emph{Say it:} \emph{udyān}
  \item \emph{Say it:} \emph{udyān}, once more
  \item \emph{Say it:} say \emph{udyān}
  \item \emph{Recall:} say the Marathi for a square, then the Marathi for a garden, then say \emph{udyān} again
\end{itemize}

\begin{tcolorbox}[breakable,colback=teal!4,colframe=teal!35!black,title={Before you move on}]
What does \mr{उद्यान} mean? (\textquotedblleft{}A park\textquotedblright{}.) Say it once more.
\end{tcolorbox}
\section[\texorpdfstring{\mr{विमानतळ} (vimāntaḷ)}{विमानतळ (vimāntaḷ)}]{\mr{विमानतळ} (vimāntaḷ) — an airport}

\label{lesson:MR-C128-vimantal}



\begin{quote}
Before the new one: say the Marathi for a garden, then the Marathi for a park.
\end{quote}

\subsection*{You'll want to know: \mr{विमानतळ}}
\textbf{\mr{विमानतळ}} — \emph{vimāntaḷ} — \textquotedblleft{}an airport\textquotedblright{}.

\begin{cousinweb}
\textbf{\mr{विमानतळ}} (\emph{vimāntaḷ}) is a modern compound of Sanskrit parts: \emph{vimāna}, \textquotedblleft{}car of the gods\textquotedblright{}, now \textquotedblleft{}aircraft\textquotedblright{} (CDIAL 11823), and \emph{taḷ}, \textquotedblleft{}surface, ground\textquotedblright{}, from Sanskrit \emph{tala} (CDIAL 5731).
\end{cousinweb}

\begin{grammarlens}[title={What you've built}]
One more word for this chapter.
\end{grammarlens}

\subsection*{Your turn}
\begin{itemize}
\raggedright
  \item \emph{Say it:} \emph{vimāntaḷ}
  \item \emph{Say it:} \emph{vimāntaḷ}, once more
  \item \emph{Say it:} say \emph{vimāntaḷ}
  \item \emph{Recall:} say the Marathi for a garden, then the Marathi for a park, then say \emph{vimāntaḷ} again
\end{itemize}

\begin{tcolorbox}[breakable,colback=teal!4,colframe=teal!35!black,title={Before you move on}]
What does \mr{विमानतळ} mean? (\textquotedblleft{}An airport\textquotedblright{}.) Say it once more.
\end{tcolorbox}
\section[\texorpdfstring{\mr{रुग्णालय} (rugṇālay)}{रुग्णालय (rugṇālay)}]{\mr{रुग्णालय} (rugṇālay) — a hospital}

\label{lesson:MR-C128-rugnalay}



\begin{quote}
Before the new one: say the Marathi for a park, then the Marathi for an airport.
\end{quote}

\subsection*{You'll want to know: \mr{रुग्णालय}}
\textbf{\mr{रुग्णालय}} — \emph{rugṇālay} — \textquotedblleft{}a hospital\textquotedblright{}.

\begin{cousinweb}
\textbf{\mr{रुग्णालय}} (\emph{rugṇālay}) is a modern compound of Sanskrit parts, a tatsama: \emph{rugṇa}, \textquotedblleft{}broken, diseased\textquotedblright{} (CDIAL 10760), and \emph{ālaya}, \textquotedblleft{}abode\textquotedblright{}: a house for the sick.
\end{cousinweb}

\begin{grammarlens}[title={What you've built}]
One more word for this chapter.
\end{grammarlens}

\subsection*{Your turn}
\begin{itemize}
\raggedright
  \item \emph{Say it:} \emph{rugṇālay}
  \item \emph{Say it:} \emph{rugṇālay}, once more
  \item \emph{Say it:} say \emph{rugṇālay}
  \item \emph{Recall:} say the Marathi for a park, then the Marathi for an airport, then say \emph{rugṇālay} again
\end{itemize}

\begin{tcolorbox}[breakable,colback=teal!4,colframe=teal!35!black,title={Before you move on}]
What does \mr{रुग्णालय} mean? (\textquotedblleft{}A hospital\textquotedblright{}.) Say it once more.
\end{tcolorbox}
